\chapter[PROFIL PERUSAHAAN]{\\ PROFIL PERUSAHAAN}
\label{chap:profil-perusahaan}

\section{Sejarah Singkat Perusahaan}
\label{sec:sejarah-singkat-perusahaan}

PT Pembangkitan Jawa Bali Masdar Solar Energi (PT PMSE) merupakan perusahaan pengelola proyek Pembangkit Listrik Tenaga Surya (PLTS) Terapung Cirata 145 MWac. Proyek ini dikembangkan untuk memanfaatkan area Waduk Cirata sebagai lokasi pembangkitan listrik berbasis energi baru dan terbarukan. Pemilihan Waduk Cirata berkaitan dengan ketersediaan area perairan yang luas, kedekatan dengan infrastruktur ketenagalistrikan, serta potensi pengembangan PLTS terapung skala besar.

PLTS Terapung Cirata berlokasi di Waduk Cirata, Jawa Barat, pada wilayah Kabupaten Bandung Barat dan Kabupaten Purwakarta. Proyek ini memiliki kapasitas 145 MWac dengan PLN sebagai \textit{off-taker}. Skema komersial yang digunakan adalah \textit{Power Purchase Agreement} (PPA) selama 25 tahun dengan lingkup pekerjaan yang meliputi pembangunan pembangkit, gardu induk, dan jalur transmisi 150 kV \cite{pmse2021}.

Pengembangan PLTS Terapung Cirata ditujukan untuk meningkatkan pemanfaatan energi surya, mendukung peningkatan bauran energi baru dan terbarukan, serta menjadi proyek percontohan PLTS terapung skala besar di Indonesia. Proyek ini juga menunjukkan adanya kerja sama pengembangan energi antara Indonesia dan Uni Emirat Arab melalui keterlibatan PT PJB dan Masdar. Visualisasi PLTS Terapung Cirata ditunjukkan pada Gambar~\ref{gbr:visualisasi-plts-cirata}.

\begin{figure}[H]
    \centering
    \includegraphics[width=0.86\textwidth]{fig_2_1_visualisasi_plts_cirata.png}
    \caption{Visualisasi PLTS Terapung Cirata pada area Waduk Cirata \cite{pmse2021}.}
    \label{gbr:visualisasi-plts-cirata}
    \vspace{-0.20cm}
\end{figure}

\section{Struktur Organisasi}
\label{sec:struktur-organisasi}

PLTS Terapung Cirata 145 MWac dimiliki dan dikelola oleh PT PMSE. PT PMSE merupakan perusahaan konsorsium atau \textit{joint venture} antara PT Pembangkitan Jawa Bali (PT PJB) dan Abu Dhabi Future Energy Company (Masdar). Komposisi kepemilikan saham terdiri atas 51\% oleh PT PJB dan 49\% oleh Masdar \cite{pmse2021}.

PT PJB merupakan anak perusahaan PT PLN (Persero), sedangkan Masdar merupakan anak usaha dari Mubadala Investment Company. Dengan susunan tersebut, proyek PLTS Terapung Cirata melibatkan pihak Indonesia melalui PLN dan PT PJB, serta pihak Uni Emirat Arab melalui Masdar. Struktur kepemilikan PT PMSE ditunjukkan pada Gambar~\ref{gbr:struktur-kepemilikan-pmse}.

\begin{figure}[H]
    \centering
    \includegraphics[width=0.62\textwidth]{fig_2_2_struktur_kepemilikan.png}
    \caption{Struktur kepemilikan PT PMSE \cite{pmse2021}.}
    \label{gbr:struktur-kepemilikan-pmse}
    \vspace{-0.20cm}
\end{figure}

Pada tingkat manajemen, struktur PT PMSE mencakup fungsi pimpinan perusahaan, operasi, serta keuangan dan administrasi. Pembagian fungsi tersebut diperlukan agar pengelolaan proyek, pelaksanaan operasi, dan administrasi perusahaan dapat berjalan secara terkoordinasi.

\section{Tugas dan Tanggung Jawab}
\label{sec:tugas-tanggung-jawab}

PT PMSE berperan sebagai pemilik dan pengelola proyek PLTS Terapung Cirata. Dalam skema PPA, tanggung jawab utama perusahaan mencakup penyediaan pembangkit, gardu induk, dan jalur transmisi 150 kV. Dengan demikian, peran PT PMSE tidak hanya terbatas pada pengelolaan pembangkit, tetapi juga mencakup penyiapan infrastruktur agar energi listrik yang dihasilkan dapat disalurkan ke sistem kelistrikan PLN.

Pelaksanaan proyek melibatkan beberapa pihak dengan lingkup pekerjaan yang berbeda. Power China Corporation Group ditunjuk sebagai kontraktor \textit{Engineering, Procurement, and Construction} (EPC) utama. PT Tractebel Indonesia berperan sebagai \textit{Project Management Consultant} (PMC) atau \textit{owner engineer}, sedangkan DNV GL berperan sebagai \textit{Lender Technical Advisor} yang memberikan laporan teknis kepada pihak pemberi pinjaman \cite{pmse2021}. Hubungan pihak-pihak tersebut ditunjukkan pada Gambar~\ref{gbr:struktur-proyek-plts-cirata}.

\begin{figure}[H]
    \centering
    \includegraphics[width=0.62\textwidth]{fig_2_3_struktur_proyek.png}
    \caption{Struktur proyek PLTS Terapung Cirata \cite{pmse2021}.}
    \label{gbr:struktur-proyek-plts-cirata}
    \vspace{-0.20cm}
\end{figure}

Berdasarkan lingkup tersebut, tugas dan tanggung jawab umum PT PMSE pada proyek PLTS Terapung Cirata dapat dirangkum sebagai berikut.

\begin{packed_enum}
    \item Mengelola pengembangan dan pengoperasian proyek PLTS Terapung Cirata.
    \item Menyediakan pembangkit, gardu induk, dan jalur transmisi 150 kV sesuai lingkup PPA.
    \item Melakukan koordinasi dengan PLN, kontraktor EPC, konsultan manajemen proyek, dan penasihat teknis lender.
    \item Menjaga keberlanjutan operasi pembangkit melalui kegiatan inspeksi, pemeliharaan, dan evaluasi teknis.
\end{packed_enum}

Secara umum, tahapan pengembangan PLTS Terapung Cirata meliputi inisiasi, studi kelayakan, pengajuan proyek, pemilihan mitra, penyusunan PPA, pembentukan perusahaan proyek, pengurusan perizinan, pembiayaan, konstruksi, hingga operasi. Rangkaian tersebut menunjukkan bahwa proyek PLTS terapung memerlukan koordinasi lintas bidang, mulai dari aspek teknis, finansial, legal, perizinan, konstruksi, sampai operasi dan pemeliharaan.

Dari sisi teknis, sistem PLTS Terapung Cirata tersusun atas modul surya, sistem pelampung, DC combiner, inverter, transformator penaik tegangan, sistem \textit{anchoring} dan \textit{mooring}, serta jalur transmisi. Komponen-komponen tersebut bekerja secara terintegrasi untuk mengubah radiasi matahari menjadi energi listrik yang dapat disalurkan ke jaringan. Susunan komponen utama PLTS terapung ditunjukkan pada Gambar~\ref{gbr:teknologi-plts-terapung}.

\begin{figure}[H]
    \centering
    \includegraphics[width=0.86\textwidth]{fig_2_4_teknologi_plts_terapung.png}
    \caption{Komponen utama teknologi PLTS terapung \cite{pmse2021}.}
    \label{gbr:teknologi-plts-terapung}
    \vspace{-0.20cm}
\end{figure}

\section{Kesehatan dan Keselamatan Kerja (K3)}
\label{sec:k3}

Kegiatan kerja praktik di lingkungan PLTS Terapung Cirata memerlukan perhatian terhadap aspek kesehatan dan keselamatan kerja karena aktivitas dilakukan pada area pembangkit, area perairan, dan fasilitas kelistrikan. Penerapan K3 dilakukan dengan mengikuti arahan pembimbing lapangan, ketentuan perusahaan, serta instruksi HSE yang berlaku di lokasi. Penggunaan alat pelindung diri, pembatasan akses ke area tertentu, dan pendampingan ketika berada di area kerja menjadi bagian penting dalam pelaksanaan kegiatan lapangan. Rincian mengenai aspek-aspek K3 yang diterapkan selama kerja praktik tercantum pada Tabel~\ref{tab:k3-kerja-praktik}.

Sebelum melaksanakan aktivitas di area \textit{offshore} PLTS Terapung Cirata, mahasiswa kerja praktik terlebih dahulu mengikuti pengenalan keselamatan kerja. Kegiatan ini membahas aturan dasar K3, potensi bahaya di area pembangkit dan perairan, penggunaan alat pelindung diri, prosedur akses area kerja, serta mekanisme komunikasi dengan petugas lapangan. Dalam pelaksanaan kerja praktik, pengenalan K3 diberikan oleh HSE HDEC/PowerChina yang berada dalam lingkup kontraktor EPC proyek. Di sisi lain, PT PMSE sebagai pengelola PLTS Terapung Cirata juga memiliki fungsi HSE untuk memastikan kegiatan operasi dan pemeliharaan berjalan sesuai prosedur keselamatan. Kegiatan pengenalan K3 tersebut ditunjukkan pada Gambar~\ref{gbr:pengenalan-k3-hdec}.

\begin{figure}[H]
    \centering
    \includegraphics[width=0.86\textwidth]{fig_2_5_pengenalan_k3_hdec.jpeg}
    \caption{Pengenalan K3 sebelum kegiatan lapangan di PLTS Terapung Cirata.}
    \label{gbr:pengenalan-k3-hdec}
    \vspace{-0.20cm}
\end{figure}

Selain pengenalan K3, mahasiswa juga mengikuti kegiatan \textit{Basic Sea Survival} (BSS) sebagai pembekalan sebelum melakukan aktivitas pada area perairan. Kegiatan BSS diberikan untuk meningkatkan kesiapan peserta dalam menghadapi kondisi darurat di air, terutama karena sebagian area kerja PLTS Terapung Cirata berada di atas Waduk Cirata. Materi dan latihan yang diberikan meliputi pengenalan alat keselamatan, penggunaan pelampung, teknik mengapung, serta prosedur dasar penyelamatan diri ketika berada di perairan. Latihan Basic Sea Survival (BSS) mahasiswa ditunjukkan pada Gambar~\ref{gbr:basic-sea-survival}.

\begin{figure}[H]
    \centering
    \includegraphics[width=0.92\textwidth]{fig_2_6_basic_sea_survival.png}
    \caption{Kegiatan \textit{Basic Sea Survival} sebelum aktivitas di area perairan.}
    \label{gbr:basic-sea-survival}
    \vspace{-0.20cm}
\end{figure}

\begin{table}[H]
    \centering
    \caption{Aspek K3 selama kerja praktik di PLTS Terapung Cirata}
    \label{tab:k3-kerja-praktik}
    \vspace{0.05cm}
    {\small
    \renewcommand{\arraystretch}{1.20}
    \begin{tabularx}{\textwidth}{
        >{\centering\arraybackslash}p{1.0cm}
        >{\raggedright\arraybackslash}p{4.0cm}
        >{\raggedright\arraybackslash}X
    }
        \toprule
        \textbf{No.} & \textbf{Nama} & \textbf{Fungsi} \\
        \midrule
        1 & \textit{Pengenalan K3} & Memberikan pemahaman mengenai aturan keselamatan, potensi bahaya, penggunaan APD, dan prosedur akses area kerja. \\
        \midrule
        2 & \textit{Basic Sea Survival} & Membekali peserta dengan kemampuan dasar keselamatan di perairan sebelum aktivitas \textit{offshore}. \\
        \midrule
        3 & \textit{Alat Pelindung Diri} & Mengurangi risiko kecelakaan melalui penggunaan helm keselamatan, sepatu keselamatan, rompi, dan pelampung sesuai kebutuhan area kerja. \\
        \midrule
        4 & \textit{Pendampingan Lapangan} & Memastikan kegiatan mahasiswa dilakukan sesuai batas kewenangan, arahan pembimbing, dan prosedur keselamatan di lokasi. \\
        \bottomrule
    \end{tabularx}}
    \vspace{-0.15cm}
\end{table}

Pada sisi operasi dan pemeliharaan, PLTS Terapung Cirata memiliki jalur \textit{gangway} sebagai akses teknisi. Akses menuju \textit{array} dilakukan ketika terdapat kegiatan inspeksi, perbaikan, penggantian komponen, atau pemeriksaan visual secara detail. Selama kerja praktik, aktivitas lapangan dilakukan dengan memperhatikan batas kewenangan mahasiswa, arahan pembimbing, dan prosedur keselamatan yang berlaku di lokasi. Akses inspeksi dan pemeliharaan pada PLTS terapung ditunjukkan pada Gambar~\ref{gbr:pemeliharaan-plts-terapung}.

\begin{figure}[H]
    \centering
    \includegraphics[width=0.78\textwidth]{fig_2_6_proses_pemeliharaan.png}
    \caption{Kegiatan inspeksi dan pemeliharaan pada PLTS terapung \cite{pmse2021}.}
    \label{gbr:pemeliharaan-plts-terapung}
    \vspace{-0.20cm}
\end{figure}

\section{Etika Profesi}
\label{sec:etika-profesi}

Etika profesi selama kerja praktik diterapkan melalui kedisiplinan, tanggung jawab, dan kepatuhan terhadap arahan pembimbing lapangan. Mahasiswa perlu menjaga perilaku profesional di lingkungan kerja, menghormati aturan perusahaan, serta menggunakan waktu kerja praktik untuk memperoleh pemahaman teknis secara aktif dan bertanggung jawab.

Mahasiswa juga perlu menjaga kerahasiaan data dan informasi yang diperoleh selama kegiatan kerja praktik, terutama informasi teknis, dokumentasi internal, dan data proyek yang tidak diperkenankan untuk disebarluaskan. Penggunaan dokumentasi, foto, maupun data pengukuran dilakukan dengan memperhatikan izin dari pihak yang berwenang.

Dalam penyusunan laporan, etika profesi diwujudkan melalui kejujuran dalam pengolahan data, ketelitian dalam penulisan, serta pencantuman sumber informasi yang digunakan. Dengan demikian, laporan kerja praktik tidak hanya menjadi dokumentasi kegiatan, tetapi juga menjadi bentuk pertanggungjawaban akademik dan profesional mahasiswa teknik elektro.
